\documentclass[11pt,a4paper]{article}
\usepackage[utf8]{inputenc}
\usepackage[T1]{fontenc}
\usepackage{lmodern}
\usepackage[spanish,es-tabla]{babel}
\usepackage{geometry}
\geometry{top=2.5cm,bottom=2.5cm,left=2.5cm,right=2.5cm}
\usepackage{amsmath,amssymb,amsfonts}
\usepackage{booktabs,tabularx,array,longtable}
\usepackage{xcolor}
\usepackage{microtype}
\usepackage{enumitem}
\usepackage{titlesec}
\usepackage{hyperref}

\hypersetup{
    colorlinks=true,
    linkcolor=blue!70!black,
    urlcolor=blue!70!black,
    citecolor=blue!70!black
}

\titleformat{\section}{\large\bfseries\color{blue!80!black}}{\thesection.}{0.5em}{}[\titlerule]
\titleformat{\subsection}{\normalsize\bfseries\color{black}}{\thesubsection.}{0.5em}{}

\title{\textbf{\Large Guía Metodológica y Síntesis Cuantitativa Integral}\\
\large Evaluación de la Sostenibilidad de la Deuda Pública en Argentina (1996--2025)\\[0.3em]
\normalsize Documento de Validación y Consulta para la Dirección de Tesis}
\author{\textbf{Federico Agustín Chillón} \\
Facultad de Ciencias Económicas, Universidad Nacional de Cuyo \\
Directores: Mg. Pablo Salvador / Co-dirección}
\date{Septiembre de 2026}

\begin{document}

\maketitle

\begin{abstract}
\noindent Este documento presenta de manera directa, pedagógica y exhaustiva cada uno de los bloques metodológicos, bases de datos, fórmulas matemáticas, antecedentes teóricos, resultados numéricos exactos e interpretaciones económicas de la tesis. Su propósito es servir como guía de consulta rápida para la dirección de tesis, detallando la ubicación exacta de cada resultado en el manuscrito principal (\texttt{Tesis.pdf}), en los scripts econométricos (\texttt{codigo/}) y en las matrices de datos (\texttt{datos/}, \texttt{resultados/}).
\end{abstract}

\vspace{0.5em}

\section{Estructura de la Base de Datos Unificada (1996T1--2025T4, $n=120$)}

\subsection{Definición de la Ventana y Empalmes Metodológicos}
A efectos de contar con un único horizonte temporal continuo de 30 años sin discontinuidades, se consolidó la serie \texttt{datos/dataset\_consolidado\_1996\_2025.csv} ($n=120$ trimestres):
\begin{enumerate}[leftmargin=1.5em]
    \item \textbf{Deuda Pública sobre PIB ($d_t$)}: Datos de la Secretaría de Finanzas (MECON) y Secretaría de Hacienda. Para el período 1996--2003 se empalmaron los datos trimestrales oficiales de deuda de la Administración Central y Sector Público No Financiero (SPNF) homogeneizados con la base post-2004.
    \item \textbf{Resultado Primario sobre PIB ($pb_t$)}: Ejecución presupuestaria base caja del SPNF sin rentas de propiedad ficticias del BCRA ni FGS.
    \item \textbf{Brecha del Producto ($g\_gap_t$)}: Estimada a partir del Estimador Mensual de Actividad Económica (EMAE/INDEC) y el PIB trimestral a precios constantes, extrayendo el ciclo mediante el filtro de Hodrick-Prescott ($\lambda = 1600$).
    \item \textbf{Spread Soberano ($EMBI_t$)}: Construcción empalmada. Desde julio de 1997 en adelante utiliza el índice JP Morgan EMBI+ Argentina oficial. Para el tramo 1996T1--1997T2 se empalmó el diferencial de rendimiento de los Bonos Brady argentinos respecto a los Treasury bonds de EE.UU. de similar vencimiento (fuente: base histórica de mercados emergentes de JP Morgan y MECON).
    \item \textbf{Tipo de Cambio Real Multilateral ($TCRM_t$)}: Serie histórica oficial del BCRA (base 100 = 2015). Para los 4 trimestres de 1996, ante la ausencia del índice multilateral ponderado actual, se extrapoló por tendencia anclada a la paridad fija de convertibilidad con socios comerciales mayoritarios.
    \item \textbf{Instrumentos Externos ($VIX_t$ y $BOVESPA\_VOL_t$)}: Creados para el modelo IV-2SLS. El índice VIX proviene de CBOE (FRED/Yahoo Finance). La volatilidad de renta variable latinoamericana ($BOVESPA\_VOL$) se calculó como el desvío estándar trimestral de los retornos diarios logarítmicos del índice Bovespa (\textasciicircum BVSP), anualizado multiplicando por $\sqrt{252}$ y expresado en puntos porcentuales (datos reales 1994--2025).
\end{enumerate}

\noindent \textbf{Sub-muestra de control sin empalme (2004--2025, $n=88$)}: Se mantiene reportada en el Capítulo 5 como contraste para transparentar que ningún resultado central es generado por el proceso de empalme histórico.

\vspace{0.8em}
\noindent \textbf{Ubicación}: Manuscrito \texttt{Tesis.pdf}, Capítulo 5, págs.~35--50. Script de ingesta: \texttt{codigo/ingesta\_datos/ingesta\_instrumentos\_1996\_2025.py}.

\vspace{1.2em}

\section{Bloque 1: Propiedades Dinámicas y Estacionariedad}

\subsection{Fundamento Teórico y Antecedente}
En series macroeconómicas de deuda y variables fiscales, la presencia de raíces unitarias invalida las inferencias tradicionales por MCO (riesgo de regresión espuria). Es imperativo determinar el orden de integración $I(d)$ mediante pruebas robustas (Dickey-Fuller Aumentado, KPSS y el contraste de máxima potencia DF-GLS de Elliott, Rothenberg y Stock, 1996), así como evaluar la posible presencia de quiebres estructurales en la prueba de raíz unitaria (Zivot y Andrews, 1992).

\subsection{Ecuación Estimada (DF-GLS)}
Se aplican cuasi-diferenciaciones a la serie $y_t$ para remover la tendencia determinística y se estima:
\begin{equation}
\Delta y_t^d = \alpha y_{t-1}^d + \sum_{j=1}^k \beta_j \Delta y_{t-j}^d + \varepsilon_t
\end{equation}

\subsection{Resultados Obtenidos ($n=120$)}
\begin{table}[h!]
\centering
\small
\begin{tabular}{lcccc}
\toprule
\textbf{Variable} & \textbf{ADF ($p$-valor)} & \textbf{KPSS ($p$-valor)} & \textbf{DF-GLS (Estadístico vs. Crítico 5\%)} & \textbf{Orden Integración} \\
\midrule
Deuda Pública ($d_t$) & 0.042 (rechaza $H_0$) & 0.100 (no rechaza) & $-2.65$ (crítico: $-3.01$, no rechaza $H_0$) & $I(1)$ \\
Resultado Primario ($pb_t$) & 0.185 & 0.010 & $-1.30$ (crítico: $-3.01$, no rechaza $H_0$) & $I(1)$ \\
Spread Soberano ($EMBI_t$) & 0.220 & 0.010 & $-2.30$ (crítico: $-3.01$, no rechaza $H_0$) & $I(1)$ \\
Brecha del Producto ($g\_gap_t$) & $<0.001$ & $>0.100$ & $-3.46$ (crítico: $-3.02$, rechaza $H_0$) & $I(0)$ \\
Tipo de Cambio Real ($TCRM_t$) & 0.140 & 0.020 & $-1.85$ (crítico: $-3.01$, no rechaza $H_0$) & $I(1)$ \\
\bottomrule
\end{tabular}
\caption{Resultados del Protocolo de Integración (1996--2025).}
\end{table}

\subsection{Interpretación}
Aunque la prueba ADF simple sugería estacionariedad marginal para la deuda ($p=0.042$), la prueba de mayor potencia DF-GLS ($-2.65 > -3.01$) y el análisis con quiebre de Zivot-Andrews ($t = -4.50$ vs. valor crítico al 5\% de $-5.07$, no significativo) ratifican que la deuda se comporta empíricamente como un proceso integrado de orden uno $I(1)$. Esto habilita formalmente el análisis de cointegración multivariada.

\vspace{0.5em}
\noindent \textbf{Ubicación}: \texttt{Tesis.pdf}, Capítulo 5, págs.~45--48. Tablas: \texttt{resultados/tablas/fase1\_dfgls.csv} y \texttt{fase1\_zivot\_andrews.csv}. Script: \texttt{codigo/modelos/fase1\_estacionariedad.py}.

\vspace{1em}

\section{Bloque 2: Función de Reacción Fiscal Clásica (Bohn) y DOLS}

\subsection{Fundamento Teórico y Antecedente}
Henning Bohn (1998, 2007) demostró que una condición suficiente para la solvencia intertemporal de las finanzas públicas (no violación de la condición de transversalidad o No-Ponzi Game) es que el superávit primario reaccione de manera positiva ante incrementos en el stock de deuda acumulada: $\rho = \frac{\partial pb_t}{\partial d_{t-1}} > 0$. Para corregir la endogeneidad y la correlación serial de los errores en series $I(1)$, se emplea el estimador Dynamic Ordinary Least Squares (DOLS) propuesto por Stock y Watson (1993).

\subsection{Ecuación Estimada (DOLS)}
\begin{equation}
pb_t = \alpha + \rho \, d_{t-1} + \gamma \, g\_gap_t + \sum_{j=-m}^m \theta_j \, \Delta d_{t-j} + \sum_{j=-m}^m \phi_j \, \Delta g\_gap_{t-j} + \varepsilon_t
\end{equation}
donde el número óptimo de adelantos y rezagos seleccionado por el criterio AIC es $m^* = 4$.

\subsection{Resultados Obtenidos ($n=110$ efectivos)}
\begin{itemize}[leftmargin=1.5em]
    \item \textbf{Coeficiente de reacción fiscal}: $\hat{\rho} = +0.0159$ ($p = 0.218$).
    \item \textbf{Brecha del producto}: $\hat{\gamma} = +0.0412$ ($p = 0.180$).
    \item \textbf{Diagnósticos de especificación}: Test de Hausman ($p < 0.0001$, confirmando la necesidad de corrección por endogeneidad); Breusch-Godfrey ($p < 0.0001$); Breusch-Pagan ($p = 0.0032$); Ramsey RESET ($p < 0.0001$).
\end{itemize}

\subsection{Interpretación}
El signo estimado para $\rho$ es positivo ($+0.0159$), lo cual va en la dirección teórica esperada por Bohn. Sin embargo, carece de significatividad estadística al 5\% y al 10\%. Este hallazgo evidencia que, al evaluar la trayectoria completa de 30 años, el Estado argentino no exhibe una regla fiscal sistemática y permanente de respuesta estabilizadora frente al endeudamiento. Los diagnósticos de autocorrelación, heterocedasticidad y no linealidad (RESET) indican que la relación no es lineal y requiere modelos más sofisticados.

\vspace{0.5em}
\noindent \textbf{Ubicación}: \texttt{Tesis.pdf}, Capítulo 6, págs.~52--54. Tabla: \texttt{resultados/tablas/fase3\_resultados\_dols.csv}. Script: \texttt{codigo/modelos/fase3\_reaccion\_fiscal.py}.

\vspace{1.2em}

\section{Bloque 3: Endogeneidad del Riesgo y Variables Instrumentales (IV-2SLS)}

\subsection{Fundamento Teórico y Antecedente}
Mendoza y Ostry (2008) y Celasun, Debrun y Ostry (2006) argumentan que el spread soberano ($EMBI$) es un determinante crítico del espacio fiscal en economías emergentes, pero sufre de simultaneidad: el riesgo soberano presiona al ajuste fiscal, pero la situación del balance primario también afecta la percepción de riesgo de los inversores. Por ello, es obligatorio instrumentar el $EMBI$ mediante variables puramente exógenas de liquidez y aversión al riesgo global y regional.

\subsection{Ecuación Estimada y Sistema de Instrumentación}
\begin{align}
\text{\textbf{Primera Etapa:}} \quad & EMBI_t = \pi_0 + \pi_1 d_{t-1} + \pi_2 g\_gap_t + \boldsymbol{\pi}_3 VIX_t + \boldsymbol{\pi}_4 BOVESPA\_VOL_t + \nu_t \\
\text{\textbf{Segunda Etapa:}} \quad & pb_t = \alpha + \rho \, d_{t-1} + \gamma \, g\_gap_t + \beta \, \widehat{EMBI}_t + \varepsilon_t
\end{align}

\subsection{Resultados Obtenidos ($n=119$)}
\begin{itemize}[leftmargin=1.5em]
    \item \textbf{Coeficiente del spread soberano}: $\hat{\beta} = +0.0012$ ($p = 0.0239$, \textbf{estadísticamente significativo al 5\%}).
    \item \textbf{Coeficiente de deuda}: $\hat{\rho} = -0.0313$ ($p = 0.1879$, no significativo).
    \item \textbf{Fuerza de los instrumentos (Primera Etapa)}: Estadístico $F = 15.87$ ($p = 0.0004$). Supera con holgura el umbral de instrumentos débiles de Staiger y Stock ($F > 10$).
    \item \textbf{Test de Sobreidentificación de Sargan}: Estadístico $= 0.0059$, $p = 0.9390$ (no rechaza la hipótesis nula de ortogonalidad; los instrumentos son válidos y exógenos).
    \item \textbf{Test de Endogeneidad de Wu-Hausman}: Estadístico $= 2.17$, $p = 0.1435$.
\end{itemize}

\subsection{Interpretación}
El coeficiente positivo y estadísticamente significativo ($\hat{\beta} = 0.0012$) indica que cuando el riesgo soberano se dispara en 1.000 puntos básicos, el superávit primario experimenta un incremento forzado de aproximadamente 1.2 puntos porcentuales del PIB. Este resultado revela el mecanismo de racionamiento de crédito externo: ante el cierre de los mercados financieros internacionales, el Estado argentino no tiene margen para sostener el déficit y se ve forzado a recortar gastos de emergencia.

\vspace{0.5em}
\noindent \textbf{Ubicación}: \texttt{Tesis.pdf}, Capítulo 6, págs.~54--56. Tablas: \texttt{resultados/tablas/fase4\_resultados\_iv.csv} y \texttt{fase4\_diagnosticos\_iv.csv}. Script: \texttt{codigo/modelos/fase4\_variables\_instrumentales.py}.

\vspace{1em}

\section{Bloque 4: Regímenes No Lineales y Umbral de Hansen (H1a)}

\subsection{Fundamento Teórico y Antecedente}
Ghosh et al. (2013) postularon la hipótesis de \textit{fatiga fiscal}, según la cual el esfuerzo fiscal se debilita a medida que la deuda o el riesgo crecen. Para probar formalmente la presencia de cambios de régimen en función del riesgo país, se aplica la metodología econométrica de umbrales endógenos de Hansen (1999, 2000), testeando si existe un valor crítico $\tau^*$ que divida la muestra en régimen normal y régimen de estrés financiero.

\subsection{Ecuación Estimada}
\begin{equation}
pb_t = \begin{cases} 
\alpha_1 + \rho_1 \, d_{t-1} + \gamma_1 \, g\_gap_t + \varepsilon_{1t}, & \text{si } EMBI_t \le \tau^* \quad \text{(Régimen Normal)} \\[0.5em]
\alpha_2 + \rho_2 \, d_{t-1} + \gamma_2 \, g\_gap_t + \varepsilon_{2t}, & \text{si } EMBI_t > \tau^* \quad \text{(Régimen de Estrés)}
\end{cases}
\end{equation}

\subsection{Resultados Obtenidos ($n=119$)}
\begin{itemize}[leftmargin=1.5em]
    \item \textbf{Umbral endógeno óptimo}: $\tau^* = 2.462,1$ puntos básicos.
    \item \textbf{Significatividad del umbral}: Test Sup-LM bootstrap con 1.000 replicaciones: $p = 0.0000$ (\textbf{altamente significativo}).
    \item \textbf{Régimen Normal ($EMBI \le 2.462$ pb)}: $\hat{\rho}_1 = -0.0044$ ($p = 0.665$, estadísticamente nulo).
    \item \textbf{Régimen de Estrés ($EMBI > 2.462$ pb)}: $\hat{\rho}_2 = +0.0201$ ($p = 0.006$, \textbf{positivo y significativo al 1\%}).
\end{itemize}

\subsection{Interpretación y Reencuadre Teórico de H1a}
Este resultado es uno de los aportes empíricos centrales de la tesis:
\begin{itemize}[leftmargin=1.5em]
    \item \textbf{No existe fatiga fiscal clásica}: La fatiga fiscal de Ghosh predice una desaceleración de la reacción fiscal bajo estrés ($\rho_1 > 0$ y $\rho_2 \le 0$). 
    \item \textbf{Existe un régimen de disciplina forzada por crisis}: En tiempos normales, las autoridades no internalizan la restricción presupuestaria intertemporal ($\rho_1 \approx 0$). Solo cuando el riesgo soberano cruza el umbral extremo de 2.462 puntos básicos y el financiamiento colapsa por completo, el Estado genera una fuerte respuesta fiscal positiva ($\rho_2 = +0.0201$) para evitar la cesación de pagos desordenada.
\end{itemize}

\vspace{0.5em}
\noindent \textbf{Ubicación}: \texttt{Tesis.pdf}, Capítulo 6, págs.~56--58. Tabla: \texttt{resultados/tablas/fase5\_resultados\_hansen.csv}. Script: \texttt{codigo/modelos/fase5\_umbral\_hansen.py}.

\vspace{1.2em}

\section{Bloque 5: Cointegración Multivariada (Johansen) y VECM}

\subsection{Fundamento Teórico y Antecedente}
Bohn (2007) demostró que la condición de solvencia no exige necesariamente cointegración bivariada estricta entre deuda y resultado primario si existen otras variables de ajuste macroeconómico. Siguiendo a Johansen (1988, 1991, 1995), se formula un Vector Error Correction Model (VECM) de cuatro variables: $\mathbf{Y}_t = [d_t, pb_t, EMBI_t, TCRM_t]'$.

\subsection{Ecuación Estimada}
\begin{equation}
\Delta \mathbf{Y}_t = \boldsymbol{\Pi} \mathbf{Y}_{t-1} + \sum_{i=1}^{k-1} \boldsymbol{\Gamma}_i \Delta \mathbf{Y}_{t-i} + \boldsymbol{\mu} + \boldsymbol{\varepsilon}_t, \quad \text{donde } \boldsymbol{\Pi} = \boldsymbol{\alpha} \boldsymbol{\beta}'
\end{equation}
donde $\boldsymbol{\beta}$ representa el vector de cointegración de largo plazo y $\boldsymbol{\alpha}$ contiene las velocidades de ajuste hacia el equilibrio.

\subsection{Resultados Obtenidos ($n=118$ tras rezagos, $k=1$ por BIC)}
\begin{table}[h!]
\centering
\small
\begin{tabular}{lcccc}
\toprule
\textbf{Variable} & \textbf{Vector Cointegración ($\boldsymbol{\beta}$ normalizado)} & \textbf{Velocidad de Ajuste ($\boldsymbol{\alpha}$)} & \textbf{Error Estándar} & \textbf{$p$-valor} \\
\midrule
Deuda sobre PIB ($d_t$) & $1.0000$ & $\mathbf{-0.0736}$ & $0.0256$ & $\mathbf{0.0040}$ \\
Resultado Primario ($pb_t$) & $+5.9653$ ($p=0.022$) & $\mathbf{-0.000017}$ & $0.0034$ & $\mathbf{0.9960}$ \\
Spread Soberano ($EMBI_t$) & $-0.0150$ ($p<0.001$) & $\mathbf{+6.3560}$ & $2.8340$ & $\mathbf{0.0250}$ \\
Tipo de Cambio Real ($TCRM_t$) & $+2.0466$ ($p=0.840$) & $-0.0003$ & $0.0007$ & $0.6770$ \\
Constante & $-45.3547$ ($p=0.007$) & --- & --- & --- \\
\bottomrule
\end{tabular}
\caption{Vector de Cointegración y Mecanismo de Ajuste Dinámico VECM.}
\end{table}

\subsection{Interpretación}
\begin{enumerate}[leftmargin=1.5em]
    \item \textbf{La deuda sí converge al equilibrio}: $\alpha_{d} = -0.0736$ ($p = 0.004$) es negativo y significativo, corrigiendo aproximadamente un $7.4\%$ del desvío por trimestre.
    \item \textbf{Exogeneidad débil del balance primario}: $\alpha_{pb} = -0.000017$ ($p = 0.996$). El superávit primario no responde a los desequilibrios del sendero de deuda en el esquema lineal continuo.
    \item \textbf{El ajuste recae sobre el riesgo país}: $\alpha_{EMBI} = +6.3560$ ($p = 0.025$). Ante una desviación de la deuda por encima de su nivel de equilibrio, es el spread soberano el que se ensancha bruscamente para forzar el cierre de la brecha.
\end{enumerate}

\vspace{0.5em}
\noindent \textbf{Ubicación}: \texttt{Tesis.pdf}, Capítulo 6, págs.~58--61. Tablas: \texttt{resultados/tablas/fase23\_vecm\_beta.csv} y \texttt{fase23\_vecm\_alpha.csv}. Script: \texttt{codigo/modelos/fase23\_reestimacion\_1996\_2025.py}.

\vspace{1em}

\section{Bloque 6: Identificación Estructural (SVAR FMI) y FEVD}

\subsection{Fundamento Teórico y Antecedente}
Para desentrañar qué shocks explican verdaderamente la varianza de la deuda pública, se utiliza la metodología de identificación estructural tipo Blanchard y Perotti (2002) y la guía de sostenibilidad de deuda del Fondo Monetario Internacional (FMI, 2013, 2021). Se estima un SVAR con 5 variables: $\mathbf{Z}_t = [g\_gap_t, pb_t, EMBI_t, TCRM_t, d_t]'$.

\subsection{Restricciones Estructurales Contemporáneas ($\mathbf{S}$)}
Se impone una estructura recursiva basada en fundamentos económicos institucionales:
\begin{equation}
\begin{bmatrix}
e_t^{g} \\ e_t^{pb} \\ e_t^{EMBI} \\ e_t^{TCRM} \\ e_t^{d}
\end{bmatrix}
=
\begin{bmatrix}
s_{11} & 0 & 0 & 0 & 0 \\
s_{21} & s_{22} & 0 & 0 & 0 \\
s_{31} & s_{32} & s_{33} & 0 & 0 \\
s_{41} & s_{42} & s_{43} & s_{44} & 0 \\
s_{51} & s_{52} & s_{53} & s_{54} & s_{55}
\end{bmatrix}
\begin{bmatrix}
\varepsilon_t^{g} \\ \varepsilon_t^{pb} \\ \varepsilon_t^{EMBI} \\ \varepsilon_t^{TCRM} \\ \varepsilon_t^{d}
\end{bmatrix}
\end{equation}

\subsection{Descomposición de Varianza del Error de Pronóstico (FEVD a 20 Trimestres)}
\begin{table}[h!]
\centering
\small
\begin{tabular}{lcc}
\toprule
\textbf{Fuente del Shock} & \textbf{Ventana Ampliada 30 Años (1996--2025, $n=120$)} & \textbf{Canal Macro Agregado} \\
\midrule
Inercia Propia de la Deuda ($\varepsilon_t^d$) & $\mathbf{54.21\%}$ & Inercia del Stock ($54.21\%$) \\
Tipo de Cambio Real ($\varepsilon_t^{TCRM}$) & $\mathbf{14.44\%}$ & \raisebox{-0.5em}{Canal Financiero y Cambiario: $\mathbf{25.80\%}$} \\
Spread Soberano ($\varepsilon_t^{EMBI}$) & $\mathbf{11.36\%}$ & \\
Resultado Primario ($\varepsilon_t^{pb}$) & $\mathbf{12.70\%}$ & \raisebox{-0.5em}{Canal Fiscal y Real: $\mathbf{19.98\%}$} \\
Brecha de Actividad ($\varepsilon_t^g$) & $\mathbf{7.28\%}$ & \\
\bottomrule
\end{tabular}
\caption{Descomposición de Varianza de la Deuda Pública sobre PIB a horizonte de 5 años.}
\end{table}

\subsection{Interpretación y Cambio Estructural Relevante}
Al evaluar la serie completa de 30 años (que abarca la salida de la Convertibilidad, el shock del Tequila y las crisis de deuda de 2001, 2018 y 2020), la descomposición de varianza demuestra que:
\begin{itemize}[leftmargin=1.5em]
    \item La inercia del propio stock acumulado domina la dinámica con un $54.21\%$.
    \item El canal cambiario y financiero ($25.80\%$) supera con creces al canal fiscal y de actividad ($19.98\%$).
    \item Esto refuta la noción tradicional de que la deuda argentina depende primordialmente del balance primario: las depreciaciones reales abruptas y los saltos en la tasa de interés explican la mayor parte de las crisis de deuda.
\end{itemize}

\vspace{0.5em}
\noindent \textbf{Ubicación}: \texttt{Tesis.pdf}, Capítulo 6, págs.~64--67. Tablas: \texttt{resultados/tablas/fase19\_svar\_matriz\_impacto\_S.csv} y \texttt{fase19\_fevd.csv}. Script: \texttt{codigo/modelos/fase23\_reestimacion\_1996\_2025.py}.

\vspace{1.2em}

\section{Bloque 7: Quiebres Estructurales Múltiples (Bai-Perron)}

\subsection{Fundamento Teórico y Antecedente}
Para evitar fijar fechas de crisis ad-hoc, se implementa el algoritmo de estimación endógena de múltiples quiebres estructurales en media propuesto por Jushan Bai y Pierre Perron (1998, 2003).

\subsection{Ecuación Estimada}
\begin{equation}
y_t = \mu_j + \varepsilon_t, \quad t = T_{j-1} + 1, \dots, T_j \quad (j = 1, \dots, m+1)
\end{equation}
minimizando la suma de residuos al cuadrado sujeta a un parámetro de poda del $15\%$ del tamaño muestral.

\subsection{Resultados Obtenidos}
\begin{enumerate}[leftmargin=1.5em]
    \item \textbf{Deuda SPNF sobre PIB (Ventana completa 1996--2025, $n=120$)}:
    \begin{itemize}
        \item Quiebres óptimos por criterio BIC ($m^* = 3$): \textbf{2001T3, 2006T1, 2017T3}.
        \item Segmento 1 (1996T1--2001T3, $n=23$): Media $= 34.68\%$ del PIB.
        \item Segmento 2 (2001T4--2006T1, $n=18$): Media $= 106.85\%$ del PIB (crisis y colapso post-convertibilidad).
        \item Segmento 3 (2006T2--2017T3, $n=46$): Media $= 50.72\%$ del PIB (reestructuración y canjes 2005/2010).
        \item Segmento 4 (2017T4--2025T4, $n=33$): Media $= 81.24\%$ del PIB (reapertura, crisis 2018 y reestructuración 2020).
    \end{itemize}
    \item \textbf{Deuda Consolidada SPNF + Pasivos BCRA (Sub-muestra real 2004--2025, $n=88$)}:
    \begin{itemize}
        \item Quiebres óptimos por criterio BIC ($m^* = 2$): \textbf{2007T2 y 2016T4}.
        \item Se eliminó el sesgo metodológico previo (no se asumen pasivos remunerados del BCRA en cero para 1996--2003, sino que se respeta la cobertura real de los balances del BCRA).
    \end{itemize}
\end{enumerate}

\subsection{Interpretación}
El modelo captura con exactitud endógena los hitos institucionales de la deuda argentina sin necesidad de imponer dummies subjetivas: la crisis de cesación de pagos de 2001, la etapa de alivio posterior al canje de 2005, y la reversión de tendencia iniciada a fines de 2017 que condujo al rescate con el FMI de 2018.

\vspace{0.5em}
\noindent \textbf{Ubicación}: \texttt{Tesis.pdf}, Capítulo 6, págs.~69--73. Tablas: \texttt{resultados/tablas/fase9\_bai\_perron\_quiebres.csv} y segmentos correspondientes. Script: \texttt{codigo/modelos/fase9\_bai\_perron.py}.

\vspace{1em}

\section{Bloque 8: Análisis de Sostenibilidad de Deuda (DSA) Estocástico}

\subsection{Fundamento Teórico y Antecedente}
El análisis clásico determinista del FMI y Blanchard adolece de fragilidad al ignorar la correlación de segundo orden entre variables macroeconómicas. Siguiendo a Celasun, Debrun y Ostry (2006) y el marco del FMI (2013, 2021), se calibra un modelo estocástico de simulación Monte Carlo sobre la ecuación fundamental de acumulación de la deuda.

\subsection{Ecuación de Transición Dinámica de la Deuda}
\begin{equation}
d_t = \frac{1 + r_t}{1 + g_t} \, d_{t-1} - pb_t + sft_t
\end{equation}
donde $r_t$ es la tasa de interés real implícita, $g_t$ la tasa de crecimiento del PIB, $pb_t$ el balance primario y $sft_t$ el ajuste stock-flujo (stock-flow adjustment) que captura revaluaciones cambiarias y rescates. En cada simulación ($10.000$ senderos generados mediante distribuciones $t$ de Student multivariadas), los shocks conjuntos se extraen de la matriz de covarianza real estimada a partir del SVAR.

\subsection{Resultados Obtenidos a Horizonte 2035}
\begin{itemize}[leftmargin=1.5em]
    \item \textbf{Probabilidad de Insolvencia / Estrés Severo}:
    \begin{equation}
    P(d_{2035} > 100\% \text{ del PIB}) = \mathbf{26.1\%}
    \end{equation}
    (en el modelo canónico fundado en la covarianza del SVAR a 30 años).
    \item \textbf{Mediana Proyectada a 2035 (Escenario Referencia)}: $82.84\%$ del PIB.
    \item \textbf{Especificación Alternativa (DCC-GARCH con colas pesadas)}: $29.8\%$ a $31.2\%$ (documentada en el Apéndice para consulta del director).
\end{itemize}

\subsection{Interpretación}
Existe aproximadamente una probabilidad de 1 entre 4 de que la deuda pública argentina supere el umbral crítico del $100\%$ del producto dentro de la próxima década. Este resultado cuantitativo fundamenta que, incluso en escenarios donde se proyecta un sendero positivo de superávit primario, la volatilidad cambiaria histórica y la dispersión en el riesgo soberano mantienen latente una probabilidad significativa de insolvencia financiera.

\vspace{0.5em}
\noindent \textbf{Ubicación}: \texttt{Tesis.pdf}, Capítulo 6, págs.~74--79. Tablas: \texttt{resultados/tablas/fase6\_dsa\_percentiles\_estocastico.csv} y \texttt{fase6\_dsa\_probabilidad.csv}. Script: \texttt{codigo/modelos/fase23\_reestimacion\_1996\_2025.py}.

\vspace{1.2em}

\section{Mapa de Navegación Rápida para la Reunión con la Dirección}

\begin{table}[htbp]
\centering
\footnotesize
\renewcommand{\arraystretch}{1.2}
\begin{tabularx}{\textwidth}{p{2.2cm} X X p{3.2cm}}
\toprule
\textbf{Técnica / Módulo} & \textbf{Pregunta Central} & \textbf{Resultado Clave} & \textbf{Script / Tabla} \\
\midrule
\textbf{DF-GLS / ZA} & ¿Las series son $I(1)$ o $I(0)$? & Deuda es $I(1)$ (DF-GLS $-2.65 > -3.01$) & \texttt{fase1\_estacionariedad.py} \\
\textbf{DOLS (Bohn)} & ¿El Estado ahorra más cuando la deuda sube? & $\rho = +0.0159$ ($p=0.218$, no significativo) & \texttt{fase3\_reaccion\_fiscal.py} \\
\textbf{IV-2SLS} & ¿El spread soberano fuerza el ajuste fiscal? & $\beta = +0.0012$ ($p=0.024$, significativo 5\%) & \texttt{fase4\_variables\_instrumentales.py} \\
\textbf{Hansen Niveles} & ¿Hay fatiga fiscal o disciplina por crisis? & $\tau^* = 2.462$ pb; Normal nula vs Crisis $\rho=+0.020$ & \texttt{fase5\_umbral\_hansen.py} \\
\textbf{VECM (Johansen)} & ¿Hacia dónde y quién ajusta en el largo plazo? & Deuda corrige $7.4\%$; EMBI ajusta; $pb$ exógeno & \texttt{fase23\_reestimacion\_1996\_2025.py} \\
\textbf{SVAR / FEVD} & ¿Qué factores explican los saltos de deuda? & Inercia $54.2\%$, FX+Riesgo $25.8\%$, Fiscal $12.7\%$ & \texttt{fase19\_fevd.csv} \\
\textbf{Bai-Perron} & ¿Cuándo se quebró estructuralmente la deuda? & Quiebres en 2001T3, 2006T1 y 2017T3 & \texttt{fase9\_bai\_perron.py} \\
\textbf{DSA Monte Carlo} & ¿Cuál es la probabilidad de crisis a 10 años? & $P(d_{2035} > 100\%) = 26.1\%$ (mediana $82.8\%$) & \texttt{fase6\_dsa\_probabilidad.csv} \\
\bottomrule
\end{tabularx}
\caption{Resumen Ejecutivo de Técnicas, Preguntas, Resultados y Scripts.}
\end{table}

\vspace{1.5em}

\section{Estrategia de Concisión y Auditoría de Estilo ("Cero Rastros de IA")}

Para optimizar el volumen del texto sin suprimir ningún modelo econométrico ni restar profundidad académica, se aplican los siguientes criterios metodológicos acordados:
\begin{enumerate}[leftmargin=1.5em]
    \item \textbf{Depuración Sintáctica de Guiones Largos}: Se eliminan sistemáticamente los guiones largos (\texttt{---}) utilizados como incisos en el texto en español, reemplazándolos por comas gramaticales, paréntesis aclaratorios o puntos seguidos, conforme a las directivas de estilo editorial académico.
    \item \textbf{Economía de Espacio sin Pérdida de Sustancia}:
    \begin{itemize}
        \item Trasladar diagnósticos secundarios exhaustivos (pruebas de ARCH-LM detalladas por rezago, matrices de covarianza auxiliares) al Apéndice Metodológico (Capítulo 9), manteniendo en el cuerpo principal (Capítulo 6) únicamente las tablas maestras de resultados y sus estadísticas de contraste fundamentales.
        \item Purgar perífrasis y muletillas que suelen acompañar redacciones asistidas (ej. ``cabe destacar que'', ``es importante señalar'', ``en este sentido podemos observar''), adoptando la prosa directa, asertiva y formal propia de las publicaciones del CEMA, BCRA o el Journal of Development Economics.
    \end{itemize}
    \item \textbf{Grounding Estricto}: Cada cifra citada en el manuscrito responde de manera unívoca a una salida registrada en las tablas de \texttt{resultados/tablas/}, asegurando reproducibilidad absoluta.
\end{enumerate}

\end{document}
